\section{The statute}
\label{sec:11.4}
The power to remove is statutory, but so are the agency that exercises it, the means it uses and the money it spends, and the body that granted each can withdraw it. Congress has done this once. The Homeland Security Act of 2002 abolished the Immigration and Naturalization Service once its functions had been transferred. It divided those functions between two bureaus and barred the executive from using its reorganization authority to recombine them \autocite[§§442, 451, 471]{hsa2002}. That was a separation of offices written into statute. The enforcement bureau became the Bureau of Immigration and Customs Enforcement under a reorganization plan that took effect on 1 March 2003 \autocite{hdoc108_32}.

So the removal counterpart of the domestic clause puts two questions to the legislature. The first is whether it has authorized this agency, by these means, at this scale. The second is what remains once the agency is gone, and who does it.\note{
Amnesty International traces ICE to the aftermath of 11 September and to “an extremely problematic security framework” adopted in the years that followed \autocite{Amnesty2026PressRelease}. It now asks Congress to defund ICE and rescind its current funding, and asks House members to cosponsor the Melt ICE Act as “a crucial first step” \autocite{AmnestyUSA2026Letter}.
}
The schedule leaves one function standing: a person with no protection claim, no fear of return and no tie to the country, which in practice means someone who has only just arrived, may be removed after an individual determination by an adjudicator outside the department that prosecutes. That function goes to the port of entry and the adjudicator, and nothing called interior enforcement receives it.

Three precedents say what kind of ending is meant. The 2002 Act ended the Immigration and Naturalization Service by transferring its enforcement half intact into a department organized around security, and the workforce carried its culture and its networks with it. That is succession, and it produced ICE.

The Stasi was dissolved within months of November 1989. No successor agency received its ninety-one thousand full-time employees, though some went into the police and the economy, and its records survived because citizens' committees occupied its offices from December 1989, and its Berlin headquarters in January 1990, before the shredding was finished \autocite{bpb2023stasi}; a statute of 1991 opened them to the people they concerned \autocite{stug1991}. That is dissolution without succession, and it worked because a legitimate successor state absorbed the people and paid them.

The third is the one this chapter has to answer. In May 2003 the Coalition Provisional Authority dissolved the Iraqi army and the security ministries, roughly four hundred thousand men, with a one-month termination payment \autocite{cpa2003order2}. Monthly stipends came only after the demonstrations of June, at one of which two former soldiers were shot dead, and after the first attacks on American troops \autocite[158]{allawi2007occupation}. The Authority's de-Baathification order purged by party rank and not by conduct, tens of thousands of people, schoolteachers among them, under a commission that became a political instrument \autocite{cpa2003order1,sissons2013bitter}. Dissolving the army erased the institution through which officers could have been investigated and left their networks intact, and those networks and their skills went first to the insurgency and then to the Islamic State, whose structure was designed by a former colonel in the intelligence service of Saddam's air defense force \autocite{reuter2015terror}. When the United States later paid roughly a hundred thousand Sunni fighters to stop, the arrangement held while it was paid; after the Iraqi government took it over in 2008 and 2009 it let payments lapse for months and arrested leaders, and by October 2010 hundreds had rejoined the insurgency, some bringing what they had learned of the American military with them \autocite{williams2010sunnis}.

An ICE workforce dissolved by a bare majority in a polarized polity, with a party promising to rehire it, is closer to Baghdad than to Berlin. Congress is not an occupier, but a workforce that reads the abolishing government as illegitimate behaves as if it were occupied, and the union that represents ICE's officers has taken a side: it endorsed the President in his first campaign, in 2016 \autocite{porter2016iceunion}. The eighth term of \S\ref{sec:11.4.3} is drawn from that record.

Two bills in the current Congress answer the first question. The Abolish ICE Act ends the agency by name \autocite{hr7123}. The Melt ICE Act works through its powers: it releases everyone held in DHS detention and bars federal funds for detention and monitoring \autocite{hr7190}. Two further drafts reach the schedule's bar on identifying people by their faces. None can pass a Congress that in June 2026 voted ICE and CBP \$70 billion more \autocite{GT2026SecureAmerica}, and since 2025 most of the pipeline's money has come by reconciliation and stays available through fiscal 2029, past the reach of next year's appropriations bill \autocite{obbba2025}. A Congress with the votes to end the agency will find the text already written. Appendix~\ref{sec:A} says what each bill does, clause by clause, and where the money sits.
